\chapter{Exact equations and potential functions}

\begin{orient}
\textbf{What this chapter is for.} An exact equation is a level-set problem wearing a differential disguise. If $F(x,y)=C$ describes a family of curves, then differentiating along any one of them gives $F_x\dd x+F_y\dd y=0$, and every exact equation is such a statement read backwards: you are handed $M\dd x+N\dd y=0$ and asked to find the $F$ it came from. The test $\pdv{M}{y}=\pdv{N}{x}$ is not a computational rule but an existence statement, the assertion that $(M,N)$ really is the gradient of something, and it works because mixed second partial derivatives commute. Recovering $F$ is then a two-step integration in which the constant of the first integration is a function of the other variable, and that single observation is what separates people who can do this from people who cannot. This chapter teaches the recovery as a mechanical procedure with the check built in, shows how to choose which of $M$ and $N$ to integrate first, handles the case where the equation is not exact until an integrating factor makes it so, and covers the two shortcuts worth having: reading off exact groups by inspection, and spotting exactness one order up in a second-order equation. It ends where the subject continues, at conservative vector fields and path independence.
\end{orient}

\section{A level set in disguise}

Recall the total differential. For a differentiable $F(x,y)$,
\[\dd F=\pdv{F}{x}\dd x+\pdv{F}{y}\dd y,\]
and the geometric content is that $\dd F$ measures the first-order change in $F$ produced by an arbitrary small step $(\dd x,\dd y)$. On a level curve $F=C$ there is no change, so a curve of the family satisfies $F_x\dd x+F_y\dd y=0$. Turn that around. Given
\[M(x,y)\dd x+N(x,y)\dd y=0,\]
the equation is called \emph{exact} when there is an $F$ with $F_x=M$ and $F_y=N$, and in that case the general solution is the implicit relation $F(x,y)=C$. No integration of the differential equation is required at all: the solution curves are the level curves of a function you construct once.

The simplest exact equations are ones nobody calls exact. A separable equation written as $f(x)\dd x-g(y)\dd y=0$ has $M_y=0$ and $N_x=0$, so it passes the test trivially, and the two-step recovery degenerates to $F=\int f\dd x$ followed by $g(y)=-\int g\dd y$, which is separation of variables verbatim. Exactness is therefore the correct generalisation of separability: it is what remains of ``integrate each side'' when the two sides refuse to come apart. Keeping that in mind stops the procedure from feeling arbitrary, because you have been doing the degenerate case since the beginning.

The test for whether such an $F$ exists is Clairaut's theorem. If $F_x=M$ and $F_y=N$ with $F$ twice continuously differentiable, then
\[\pdv{M}{y}=\frac{\partial^{2}F}{\partial y\,\partial x}=\frac{\partial^{2}F}{\partial x\,\partial y}=\pdv{N}{x},\]
so $M_y=N_x$ is necessary. It is also sufficient, provided $M$ and $N$ are continuously differentiable on a region with no holes in it, and the hole condition is not a technicality invented by pedants: on the punctured plane the pair $M=-y/(x^{2}+y^{2})$, $N=x/(x^{2}+y^{2})$ passes the test everywhere and has no single-valued potential, because the candidate is $\arctan(y/x)$, which gains $2\pi$ on a circuit around the origin. On any disc avoiding the origin it is exact; on the annulus it is not. For the equations in this chapter the region is a rectangle or a half-plane and the test is decisive.

\begin{keynote}[The picture]
Plot the level curves of $F$. At each point the gradient $(F_x,F_y)=(M,N)$ points across them, in the direction of steepest increase, and the solution curve is the one running perpendicular to it. The differential equation $y'=-M/N$ is just that statement written as a slope: the tangent direction $(\dd x,\dd y)$ must have zero dot product with $(M,N)$, that is $M\dd x+N\dd y=0$. Everything that looks mysterious in the algebra is obvious in this picture, including why $N=0$ produces vertical tangents (the gradient is horizontal there) and why the solution is a family rather than a single curve (one curve per value of $C$).
\end{keynote}

\begin{keynote}[What the test is really saying]
$M_y=N_x$ is the statement that the vector field $(M,N)$ is a gradient. Every consequence follows from that single reading. The solution is a level curve because gradients are perpendicular to level curves. The recovery works because a gradient is determined by its potential and vice versa. The failure of the test means the field has circulation, and that is why an integrating factor can rescue the situation: multiplying by $\mu$ rescales the field pointwise, which changes its circulation without changing its direction, hence without changing the solution curves.
\end{keynote}

That last sentence deserves emphasis, because it is the reason the two chapters on this material belong together. Multiplying $M\dd x+N\dd y=0$ by any nonvanishing $\mu(x,y)$ leaves the solution set untouched: the equation $\mu M\dd x+\mu N\dd y=0$ has exactly the same curves. So exactness is a property of the \emph{representation}, not of the equation, and every first-order equation is exact in some representation. Finding the right one is the integrating-factor problem; recovering $F$ once the representation is right is this chapter.

\tech[Recovering the potential]{Recovering the potential $F$ from $M$ and $N$}{core}
\cue The equation is in the form $M\dd x+N\dd y=0$ and the test $M_y=N_x$ has been confirmed by actual differentiation, not by appearance.
\method Two integrations with a matching step between them.
\begin{enumerate}[label=\arabic*.]
\item Integrate $M$ with respect to $x$, holding $y$ fixed: $F=\int M\dd x+g(y)$. The unknown is a \emph{function of the other variable}, never a bare constant, because anything depending on $y$ alone has zero $x$-derivative.
\item Differentiate that expression with respect to $y$ and set it equal to $N$. Every term containing $x$ must cancel; what remains is an equation $g'(y)=\bigl(\text{function of }y\text{ alone}\bigr)$.
\item Integrate $g'$ and write $F(x,y)=C$.
\end{enumerate}
The cancellation in step 2 is the built-in check. If an $x$ survives on the right, either the equation was not exact or a derivative was computed wrongly, and you have caught the error before doing any further work.

\begin{worked}
\[(2x+y^{2})\dd x+(2xy+3y^{2})\dd y=0.\]
Test: $M_y=2y$ and $N_x=2y$, so the equation is exact. Integrating $M$ in $x$,
\[F=x^{2}+xy^{2}+g(y).\]
Differentiating in $y$ and matching to $N$: $2xy+g'(y)=2xy+3y^{2}$. The $2xy$ cancels, as it must, leaving $g'=3y^{2}$ and $g=y^{3}$. Hence
\[x^{2}+xy^{2}+y^{3}=C.\]
Through $(1,1)$ the constant is $3$; following the curve to $x=1.55$ gives $y\approx0.535286$, and $x^{2}+xy^{2}+y^{3}=3.000000$ there. Note that the curve has a vertical tangent where $N=y(2x+3y)$ vanishes, at $y=0$ and $x=\sqrt3$: the level curve continues past that point but $y$ stops being a function of $x$, which is exactly why the implicit form is the honest answer.
\end{worked}

\begin{worked}[With a transcendental potential]
\[\bigl(y\cos x+2x\eu^{y}\bigr)\dd x+\bigl(\sin x+x^{2}\eu^{y}-1\bigr)\dd y=0.\]
Both cross-derivatives equal $\cos x+2x\eu^{y}$, so the equation is exact. Integrating $M$ in $x$ gives $F=y\sin x+x^{2}\eu^{y}+g(y)$, and then
\[F_y=\sin x+x^{2}\eu^{y}+g'(y)=\sin x+x^{2}\eu^{y}-1\ \Longrightarrow\ g'=-1,\ g=-y.\]
So $y\sin x+x^{2}\eu^{y}-y=C$. The two $x$-dependent terms cancelled exactly, which is the confirmation that the first integration was right.
\end{worked}

\begin{trap}
Writing $+C$ instead of $+g(y)$ after the first integration. This is the single defining error of the technique, and it is silent: it throws away every term of the potential that depends on $y$ alone, so in the first example above it would return $x^{2}+xy^{2}=C$, which satisfies nothing. The tell is that step 2 then produces an inconsistency you have no way to absorb. A related error is integrating $M$ with respect to $x$ but treating $y$ as though it were a function of $x$; it is not, during this step it is a constant.
\end{trap}

\begin{seealsob}Integrate whichever variable is easier; grouping and solving by inspection; integrating factors depending on one variable; the total differential.\end{seealsob}

Step 1 chose $M$ arbitrarily, and that choice is free. The procedure is symmetric in the two variables, and half the labour in a badly organised solution comes from integrating the wrong one. Which is wrong is usually obvious at a glance: one of $M$ and $N$ is a polynomial and the other contains a product of transcendental functions produced by the product rule, and it is the polynomial one you want.

\tech{Integrate whichever variable is easier}{medium}
\cue One of $M$, $N$ is markedly uglier than the other, typically because the potential contains a product such as $\eu^{xy}\cos2x$ whose $x$-derivative has two terms and whose $y$-derivative has one.
\method Start from $F=\int N\dd y+h(x)$ instead. Now the unknown is $h(x)$, and you match $F_x$ against $M$. The two routes produce the same $F$, so choose the one that leaves you an elementary single-variable integral. As a rule of thumb, whichever of $M$ and $N$ has fewer terms was produced by differentiating in the variable you should integrate in.

\begin{worked}
\[2xy^{3}\dd x+(3x^{2}y^{2}+4y)\dd y=0.\]
Exact, since both cross-derivatives are $6xy^{2}$. Integrating $N$ in $y$ is immediate:
\[F=x^{2}y^{3}+2y^{2}+h(x),\qquad F_x=2xy^{3}+h'(x)=2xy^{3}\ \Longrightarrow\ h'=0.\]
So $x^{2}y^{3}+2y^{2}=C$, in two lines. Going the other way costs the same here, but the next example does not have that symmetry.
\end{worked}

\begin{worked}[Where the choice matters]
\[\bigl(y\eu^{xy}\cos2x-2\eu^{xy}\sin2x+2x\bigr)\dd x+\bigl(x\eu^{xy}\cos2x-3\bigr)\dd y=0.\]
$M$ has three terms and one of them is a product-rule leftover; $N$ has two and integrates in $y$ on sight, since $\int x\eu^{xy}\cos2x\dd y=\eu^{xy}\cos2x$ with $x$ held fixed. So
\[F=\eu^{xy}\cos2x-3y+h(x).\]
Differentiating in $x$: $F_x=y\eu^{xy}\cos2x-2\eu^{xy}\sin2x+h'(x)$, and matching to $M$ gives $h'=2x$, so $h=x^{2}$ and
\[\eu^{xy}\cos2x-3y+x^{2}=C.\]
Both product-rule terms cancelled at once. Starting from $M$ instead would have required integrating $y\eu^{xy}\cos2x-2\eu^{xy}\sin2x$ with respect to $x$, which is only doable by recognising it as a derivative, that is, by already knowing the answer.
\end{worked}

\begin{trap}
Carrying out both routes and then adding the two expressions for $F$, which doubles the potential. They are two derivations of the same function, not two halves of it. If the two routes give visibly different answers, one of them contains an arithmetic error, or the two constants differ, which is harmless: $F=C$ and $F+5=C$ describe the same family.
\end{trap}

\begin{seealsob}Recovering the potential $F$ from $M$ and $N$; grouping and solving by inspection.\end{seealsob}

\section{When the equation is not exact}

Most equations that arrive in the form $M\dd x+N\dd y=0$ fail the test, and failing it is not a verdict on the equation but on the way it has been written. Because multiplying by a nonvanishing $\mu$ leaves the solution curves alone, the question is whether some $\mu$ makes the representation exact, and the systematic search for one is the business of the previous chapter. Here it is worth seeing the whole cycle run once, from a failed test through the factor to the potential, because in practice the two chapters are used as a single procedure and the seam between them should be invisible.

\begin{worked}[Not exact, then exact]
\[(3xy+y^{2})\dd x+(x^{2}+xy)\dd y=0.\]
First the test: $M_y=3x+2y$ and $N_x=2x+y$, so $M_y-N_x=x+y$ and the equation is not exact. Try the $x$-only recipe, which asks whether $(M_y-N_x)/N$ is free of $y$:
\[\frac{M_y-N_x}{N}=\frac{x+y}{x^{2}+xy}=\frac{x+y}{x(x+y)}=\frac{1}{x},\]
which it is, once the common factor $x+y$ is cancelled. Hence $\mu=\exp\bigl(\int\dd x/x\bigr)=x$, and the rescaled equation is
\[(3x^{2}y+xy^{2})\dd x+(x^{3}+x^{2}y)\dd y=0.\]
Verify before proceeding: $\partial_y(3x^{2}y+xy^{2})=3x^{2}+2xy$ and $\partial_x(x^{3}+x^{2}y)=3x^{2}+2xy$. Exact. Now integrate $M$ in $x$:
\[F=x^{3}y+\tfrac12x^{2}y^{2}+g(y),\qquad F_y=x^{3}+x^{2}y+g'(y)=x^{3}+x^{2}y,\]
so $g'=0$ and the solution is $x^{3}y+\tfrac12x^{2}y^{2}=C$, or more cleanly
\[2x^{3}y+x^{2}y^{2}=C.\]
Through $(1,1)$ the constant is $3$. Following the curve to $x=1.4$ gives $y\approx0.468318$, and $2x^{3}y+x^{2}y^{2}=3.000000$ there; at $x=2$ the relation $4y^{2}+16y=3$ gives $y\approx0.179449$. Note that the factor $\mu=x$ vanishes at $x=0$, so the derivation is valid on $x>0$ and on $x<0$ separately, and the line $x=0$ has to be examined on its own: it is itself a solution of the original equation, since $M$ and $N$ both vanish along $x=0$ only where $y=0$, and the curve $x=0$ satisfies $\dd x=0$ with $N\dd y=0$.
\end{worked}

\begin{worked}[The same cycle with a $y$-only factor]
\[2xy\dd x+\Bigl(3x^{2}+\frac{1}{y}\Bigr)\dd y=0.\]
Test: $M_y=2x$, $N_x=6x$, not exact. The $x$-only quotient is $\dfrac{M_y-N_x}{N}=\dfrac{-4x}{3x^{2}+1/y}$, which carries both variables, so try the other:
\[\frac{N_x-M_y}{M}=\frac{6x-2x}{2xy}=\frac{2}{y},\]
free of $x$. Hence $\mu=\exp\bigl(\int2\dd y/y\bigr)=y^{2}$, and the rescaled equation $2xy^{3}\dd x+(3x^{2}y^{2}+y)\dd y=0$ is exact, both cross-derivatives being $6xy^{2}$. Integrating $M$ in $x$: $F=x^{2}y^{3}+g(y)$, then $F_y=3x^{2}y^{2}+g'=3x^{2}y^{2}+y$ forces $g=\tfrac12y^{2}$, so
\[x^{2}y^{3}+\tfrac12y^{2}=C.\]
Through $(1,1)$ the constant is $\tfrac32$; at $x=1.35$ the trajectory gives $y\approx0.854046$ and $F=1.500000$. The two examples in this section are the same procedure with the two recipes, and the only judgement involved was deciding which quotient to compute first.
\end{worked}

The cancellation of $x+y$ in that computation was not luck. The recipes for one-variable integrating factors work exactly when a common factor cancels out of the quotient, so when you compute $(M_y-N_x)/N$ the first thing to do is factor the numerator and the denominator rather than expand them. An expanded quotient looks hopeless and a factored one is often a monomial.

\section{Reading the exact structure by eye}

The systematic procedure always works, and it is often slower than looking. Certain combinations of differentials are exact and recognisable on sight, and a solver who knows them can integrate an equation in one line that the procedure takes half a page to grind out. The list is short: it is essentially the derivative table of the product and quotient rules read backwards, and it is worth having in permanent memory.
\begin{center}
\footnotesize
\setlength{\tabcolsep}{5pt}
\renewcommand{\arraystretch}{1.34}
\begin{tabular}{@{}lll@{}}
\toprule
\mh{Group} & \mh{Is $\dd$ of} & \mh{Look for it when} \\
\midrule
$x\dd y+y\dd x$ & $xy$ & $M$ and $N$ are $y$ and $x$ times the same thing \\
$x\dd x+y\dd y$ & $\tfrac12\bigl(x^{2}+y^{2}\bigr)$ & $M$ and $N$ are $x$ and $y$ times the same thing \\
$\dfrac{x\dd y-y\dd x}{x^{2}}$ & $\dfrac{y}{x}$ & the equation is homogeneous of degree zero \\
$\dfrac{y\dd x-x\dd y}{y^{2}}$ & $\dfrac{x}{y}$ & the same, with the roles exchanged \\
$\dfrac{x\dd y-y\dd x}{xy}$ & $\ln\dfrac{y}{x}$ & every term carries a factor $xy$ \\
$\dfrac{x\dd y-y\dd x}{x^{2}+y^{2}}$ & $\arctan\dfrac{y}{x}$ & $x^{2}+y^{2}$ appears, or should \\
$\dfrac{x\dd x+y\dd y}{x^{2}+y^{2}}$ & $\tfrac12\ln\bigl(x^{2}+y^{2}\bigr)$ & the same, symmetric numerator \\
$\dfrac{x\dd x+y\dd y}{\sqrt{x^{2}+y^{2}}}$ & $\sqrt{x^{2}+y^{2}}$ & a radial distance is the natural variable \\
\bottomrule
\end{tabular}
\end{center}
Every entry is the product or quotient rule read backwards, and the four with an antisymmetric numerator are all the same identity divided by different things: $x\dd y-y\dd x$ is exact after division by $x^{2}$, by $y^{2}$, by $xy$, or by $x^{2}+y^{2}$, and which divisor to use is decided by what else is in the equation.
The skill is not remembering the list but seeing where to split the given expression so that one of these appears, and the split is usually forced by the antisymmetric combination $x\dd y-y\dd x$ or the symmetric one $x\dd x+y\dd y$.

\tech{Grouping and solving by inspection}{hard}
\cue $M\dd x+N\dd y$ visibly contains one of the recognisable exact blocks, most often $x\dd y+y\dd x$ or $x\dd y-y\dd x$, with the rest of the expression a pure differential in one variable.
\method Split the expression into exact groups plus a single-variable remainder, integrate each group by sight, and add. When an antisymmetric block appears, look for the divisor that turns it into an exact form: $x^{2}$, $y^{2}$, $xy$, or $x^{2}+y^{2}$, and be prepared to multiply the whole equation by that divisor's reciprocal, which is an integrating factor found by inspection rather than by recipe.

\begin{worked}
\[(x^{2}+y^{2}+y)\dd x-x\dd y=0.\]
Regroup so the antisymmetric block stands alone: $(x^{2}+y^{2})\dd x+(y\dd x-x\dd y)=0$. Divide by $x^{2}+y^{2}$, which is the integrating factor the block is asking for:
\[\dd x-\frac{x\dd y-y\dd x}{x^{2}+y^{2}}=0\ \Longrightarrow\ \dd\Bigl(x-\arctan\frac{y}{x}\Bigr)=0,\]
so the solution is $x-\arctan\dfrac{y}{x}=C$. Through $(1,1)$ the constant is $1-\tfrac{\pi}{4}\approx0.214602$, and following the trajectory to $x=1.7$ gives $y\approx19.858$, at which point $x-\arctan(y/x)=0.214602$ as required. Neither one-variable recipe would have found this factor: $(M_y-N_x)/N=-(2y+2)/x$ involves both variables, and so does the other quotient.
\end{worked}

\begin{worked}[The symmetric block, with a chain rule on top]
\[\bigl(y\cos(xy)+2x\bigr)\dd x+x\cos(xy)\dd y=0.\]
The two terms carrying the cosine are $\cos(xy)\bigl(y\dd x+x\dd y\bigr)$, and $y\dd x+x\dd y=\dd(xy)$, so that pair is $\cos(u)\dd u$ with $u=xy$, which is $\dd\bigl(\sin(xy)\bigr)$. The remaining term is $2x\dd x=\dd(x^{2})$. Hence the whole expression is $\dd\bigl(\sin(xy)+x^{2}\bigr)$ and the solution is
\[\sin(xy)+x^{2}=C.\]
Confirm by differentiating: $F_x=y\cos(xy)+2x$ and $F_y=x\cos(xy)$, both correct. The systematic method reaches the same place, but only after you notice that $\int\bigl(y\cos(xy)+2x\bigr)\dd x$ requires the substitution $u=xy$ with $y$ held fixed; seeing the group first saves that step entirely.
\end{worked}

\begin{worked}[A conserved combination, and a decorative evaluation point]
\[\dv{y}{x}=\frac{2x^{3}y-y^{4}}{x^{4}-2xy^{3}},\qquad y(1)=2.\]
Cross-multiplying, $y(2x^{3}-y^{3})\dd x=x(x^{3}-2y^{3})\dd y$. The equation is homogeneous, and the grouping route is faster: the combination $x^{3}+y^{3}$ against $xy$ is conserved, and one checks directly that $F=\dfrac{x^{3}+y^{3}}{xy}$ has $\dv{y}{x}=-F_x/F_y$ equal to the right-hand side. So the first integral is
\[x^{3}+y^{3}=Cxy,\]
and $y(1)=2$ gives $C=\tfrac92$. If the problem then asks for the value of $\dfrac{x^{3}+y^{3}}{xy}$ at some ostentatious point on the trajectory, at $x=2026$ say, the answer is $\tfrac92$ and the evaluation point is a red herring: the requested quantity is the conserved combination itself. Numerically, integrating from $(1,2)$ to $x=1.6$ gives $y\approx2.333366$ and $(x^{3}+y^{3})/(xy)=4.500000$.
\end{worked}

\begin{trap}
Being distracted by a large or ugly evaluation point. When what is asked for is the conserved combination, the point never matters; when it is $y$ itself, you must actually solve the implicit relation. Read the question before computing. The second trap is dividing by a group such as $x^{2}+y^{2}$ without noting that the division is illegal at the origin, so the resulting potential describes the solution only on a region avoiding it.
\end{trap}

\begin{seealsob}Recovering the potential $F$ from $M$ and $N$; integrating factors $x^{a}y^{b}$ and $\mu(xy)$; the total differential.\end{seealsob}

\section{Exactness one order up}

The idea generalises upward. A second-order equation is exact when one whole side of it is $\dvop{x}$ of a first-order expression, and integrating once then reduces the order by one and hands the problem to the techniques of this and the previous chapter. The pattern to look for is always the product rule or the chain rule run backwards, and the search is usually a search for the multiplier that makes it appear.

\tech{Exact second-order equations and integrable combinations}{hard}
\cue A second-order equation in which the terms could be a single derivative, most often after multiplying by a factor: the shapes $\mu y''+\mu'y'=(\mu y')'$ and $\eu^{y}y''+\eu^{y}(y')^{2}=\bigl(\eu^{y}y'\bigr)'$ are the two that recur.
\method Identify the combination and integrate once, carrying the constant. The result is a first-order equation, and the initial conditions of the original problem determine the first constant before you go any further. Then solve the first-order equation by whatever this chapter or the last one provides. The useful reflex: a term $(y')^{2}$ alongside $y''$ almost always signals that a factor of $\eu^{y}$ or $y$ is wanted, because those are the multipliers whose product rule generates a squared derivative.

\begin{worked}
\[y''+(y')^{2}=2\eu^{-y},\qquad y(0)=y'(0)=0.\]
The $(y')^{2}$ next to $y''$ is the signal. Multiply by $\eu^{y}$:
\[\eu^{y}y''+\eu^{y}(y')^{2}=2\ \Longrightarrow\ \bigl(\eu^{y}y'\bigr)'=2\ \Longrightarrow\ \eu^{y}y'=2x+C_1.\]
Both conditions are at $x=0$, where $\eu^{0}\cdot0=0$, so $C_1=0$. Now notice that $\eu^{y}y'=\bigl(\eu^{y}\bigr)'$, so $\bigl(\eu^{y}\bigr)'=2x$ and $\eu^{y}=x^{2}+C_2$, with $C_2=1$ from $y(0)=0$. Hence
\[y=\ln(x^{2}+1),\qquad y(2)=\ln5\approx1.609438,\]
confirmed to ten decimal places by numerical integration of the second-order system.
\end{worked}

\begin{worked}[The other standard multiplier]
\[yy''+(y')^{2}=0,\qquad y(0)=1,\ y'(0)=2.\]
No multiplication is needed: the left side is already $(yy')'$. So $yy'=C_1=2$, hence $\bigl(\tfrac12y^{2}\bigr)'=2$ and $y^{2}=4x+D$, with $D=1$. Thus $y=\sqrt{4x+1}$ and $y(1)=\sqrt5\approx2.236068$. The solution ceases to exist at $x=-\tfrac14$, where $y$ reaches zero and $y'$ blows up.
\end{worked}

\begin{trap}
Integrating once and then dropping the first constant before applying the initial condition on $y'$. Both constants must be carried, and the first one is determined by $y'(x_0)$ together with $y(x_0)$, not by $y(x_0)$ alone. In the first example above it is $y'(0)=0$ that kills $C_1$; had $y'(0)$ been nonzero the subsequent equation would have been $\bigl(\eu^{y}\bigr)'=2x+C_1$ and the answer a different function entirely.
\end{trap}

\begin{seealsob}Recovering the potential $F$ from $M$ and $N$; grouping and solving by inspection; reduction of order for equations with a missing variable.\end{seealsob}

\section{Reading the answer}

An exact equation returns its solution as an implicit relation $F(x,y)=C$, and that is the answer, not a step on the way to one. Three questions then arise, and they are worth separating.

The first is whether to solve for $y$. Do it only when the relation is a polynomial of degree one or two in $y$, or otherwise obviously invertible, and only when the problem asks for a value of $y$ rather than for the family of curves. Forcing an explicit form out of $x^{2}+xy^{2}+y^{3}=C$ means solving a cubic and produces three branches, of which the initial condition selects one, and nothing about the geometry is clearer afterwards.

The second is which branch. When you do solve, a sign or a root choice appears, and the initial condition fixes it: substitute the initial point and keep the branch that reproduces it. The third is the interval of validity, and it is decided by where the implicit relation stops defining $y$ as a function of $x$, which by the implicit function theorem is where $F_y=N$ vanishes. That is the same curve $N=0$ on which the solution has a vertical tangent, and it is why the systematic method never needs to divide by anything.

\begin{worked}[Implicit to explicit, with the branch chosen]
The non-exact example above gave $2x^{3}y+x^{2}y^{2}=C$, and through $(1,1)$ the constant is $3$. This is a quadratic in $y$:
\[x^{2}y^{2}+2x^{3}y-3=0,\qquad y=\frac{-2x^{3}\pm\sqrt{4x^{6}+12x^{2}}}{2x^{2}}=\frac{-x^{3}\pm\sqrt{x^{6}+3x^{2}}}{x^{2}}.\]
At $x=1$ the two branches are $-1+2=1$ and $-1-2=-3$, so the initial condition $y(1)=1$ selects the upper sign. Hence
\[y=\frac{\sqrt{x^{6}+3x^{2}}-x^{3}}{x^{2}}\qquad\text{on }x>0,\]
and $y(2)=\tfrac14\bigl(\sqrt{76}-8\bigr)\approx0.179449$, matching the numerical trajectory. The radicand is positive for every $x\neq0$, so the only obstruction is $x=0$, which is where the integrating factor $\mu=x$ died. The two facts are the same fact.
\end{worked}

\section{Choosing the route}

Two derivatives decide everything, so compute them first and think afterwards.

\begin{center}
\begin{tikzpicture}[node distance=6mm and 8mm]
\node[dtest] (a) {Compute $M_y$ and $N_x$.\\ Are they equal?};
\node[dtest, below left=9mm and 4mm of a] (b) {Which of $M$, $N$\\ has fewer terms?};
\node[dtest, below right=9mm and 2mm of a] (c) {Is a recognisable\\ exact group present?};
\node[dleaf, below left=9mm and 0mm of b] (d) {$M$: $F=\int M\dd x+g(y)$,\\ match $F_y$ to $N$};
\node[dleaf, below right=9mm and 0mm of b] (e) {$N$: $F=\int N\dd y+h(x)$,\\ match $F_x$ to $M$};
\node[dleaf, below left=9mm and 0mm of c] (f) {Divide by $x^{2}$, $xy$\\ or $x^{2}+y^{2}$; integrate by sight};
\node[dleaf, below right=9mm and 0mm of c] (g) {Integrating factor: test\\ $\dfrac{M_y-N_x}{N}$, then $\dfrac{N_x-M_y}{M}$};
\draw[dedge] (a) -- node[dlab,left]{yes} (b);
\draw[dedge] (a) -- node[dlab,right]{no} (c);
\draw[dedge] (b) -- node[dlab,left]{$M$} (d);
\draw[dedge] (b) -- node[dlab,right]{$N$} (e);
\draw[dedge] (c) -- node[dlab,left]{yes} (f);
\draw[dedge] (c) -- node[dlab,right]{no} (g);
\end{tikzpicture}
\end{center}

The right-hand branch hands the problem back to the integrating-factor chapter and returns it exact, at which point you re-enter the tree at the top. That loop is the whole algorithm, and it terminates because the list of factor shapes to try is finite.

\section{Where this goes next}

Everything in this chapter is one-dimensional shadow of a statement about vector fields, and naming the statement makes the whole subject easier to remember. The pair $(M,N)$ is a vector field on a region of the plane; the equation $M\dd x+N\dd y=0$ says that the solution curve is everywhere orthogonal to it. The field is called \emph{conservative} when it is the gradient of a scalar potential $F$, and three conditions are then equivalent on a region without holes: the field is a gradient, the test $M_y=N_x$ holds, and the line integral $\int_\Gamma M\dd x+N\dd y$ depends only on the endpoints of $\Gamma$ and not on the path. Exactness of a differential equation is the first of these; the test you apply is the second; and the third is what makes the recovery procedure work, since
\[F(x,y)=\int_{(x_0,y_0)}^{(x,y)}M\dd x+N\dd y\]
is well defined precisely because the path does not matter. The two-step integration is that line integral evaluated along a staircase path: first horizontally, which contributes $\int M\dd x$, then vertically, which contributes the $g(y)$ you solve for. Seen that way, the ``function of the other variable'' stops being a rule to remember and becomes the second leg of a journey.

\begin{worked}[The recovery, run as a line integral]
Take the first exact equation of this chapter, $(2x+y^{2})\dd x+(2xy+3y^{2})\dd y=0$, and build $F$ from the origin along a staircase. Horizontal leg, from $(0,0)$ to $(x,0)$ with $\dd y=0$ and $y=0$ throughout:
\[\int_0^{x}\bigl(2t+0\bigr)\dd t=x^{2}.\]
Vertical leg, from $(x,0)$ to $(x,y)$ with $\dd x=0$ and $x$ fixed:
\[\int_0^{y}\bigl(2xs+3s^{2}\bigr)\dd s=xy^{2}+y^{3}.\]
Adding, $F=x^{2}+xy^{2}+y^{3}$, which is exactly what the two-step procedure produced earlier, term for term. The first leg is $\int M\dd x$; the second leg is the $g(y)$ that the procedure solves for. Choosing the other staircase, vertical then horizontal, gives $\int_0^{y}3s^{2}\dd s=y^{3}$ followed by $\int_0^{x}(2t+y^{2})\dd t=x^{2}+xy^{2}$, and the same $F$: that agreement is path independence, demonstrated rather than asserted.
\end{worked}

The punctured-plane example given earlier is the standard counterexample to dropping the topological hypothesis, and it is the same object in both languages: a field with zero curl and nonzero circulation, or equivalently a differential equation that is exact locally and has no global potential. It is also where the subject stops being calculus and becomes topology, which is a good place to leave it.

The physical reading is worth one sentence too, because it explains the vocabulary. If $(M,N)$ is a force field, conservative means that the work done moving between two points does not depend on the route, so a potential energy can be defined and total energy is conserved along trajectories. The first integral of an exact second-order equation is precisely a conservation law of that kind: $yy'+\ldots$ constant is a statement that something does not change, and the ``constant of integration'' is the value of the conserved quantity fixed by the initial data. Once you read $F=C$ as a conservation law rather than as an answer with a leftover letter, the technique in this chapter and the first-integral technique for second-order equations stop being two things.

\begin{keynote}[Three checks, in increasing cost]
Before trusting a potential, run these in order. First, does $F_x=M$? One derivative. Second, does $F_y=N$? One more. If both hold, $F$ is a potential and the solution is right, whatever route produced it, so a lucky guess verified this way is as good as a derivation. Third, and only if the first two are awkward to do by hand, integrate the equation numerically from a point on the curve and confirm that $F$ has the same value at the far end. The third check catches nothing the first two miss, but it catches transcription errors, and it is the only one available when $F$ is messy enough that differentiating it is itself error-prone.
\end{keynote}

\section{Recognition matrix}
\begin{recog}{@{}p{0.30\textwidth}p{0.36\textwidth}p{0.28\textwidth}@{}}
\rowcolor{mist}\mh{If you see} & \mh{Reach for} & \mh{If that fails} \\
\midrule
\endhead
$M\dd x+N\dd y=0$, unknown type & Compute $M_y$ and $N_x$ before anything else & Rewrite $y'=f(x,y)$ as $M\dd x+N\dd y$ \\
$M_y=N_x$ confirmed & $F=\int M\dd x+g(y)$, then match $F_y$ to $N$ & Integrate $N$ in $y$ instead \\
$M$ ugly, $N$ short & Start from $F=\int N\dd y+h(x)$ & Whichever has fewer terms was differentiated \\
$x$ surviving in $g'(y)$ & The equation was not exact, or $M_y\neq N_x$ & Recheck the first integration \\
$M_y\neq N_x$ & Integrating factor: try $(M_y-N_x)/N$ first & Factor before expanding the quotient \\
A common factor in $M_y-N_x$ and $N$ & Cancel it; the quotient is usually a monomial & Then $\mu=\exp\int$ of that \\
$x\dd y+y\dd x$ present & $\dd(xy)$; group and integrate by sight & Systematic two-step recovery \\
$x\dd y-y\dd x$ present & Divide by $x^{2}$, $xy$, or $x^{2}+y^{2}$ & $\dd(y/x)$, $\dd\ln(y/x)$, $\dd\arctan(y/x)$ \\
$x\dd x+y\dd y$ present & $\tfrac12\dd\bigl(x^{2}+y^{2}\bigr)$ & Polar coordinates \\
A homogeneous right-hand side & Look for a conserved ratio such as $\tfrac{x^{3}+y^{3}}{xy}$ & Substitute $y=vx$ \\
A question asked at an absurd point & Report the conserved combination itself & Solve for $y$ only if $y$ is asked for \\
$y''$ with $(y')^{2}$ beside it & Multiply by $\eu^{y}$ or by $y$; look for $(\mu y')'$ & Reduction of order with $p=y'$ \\
$\mu y''+\mu'y'$ & $(\mu y')'$; integrate once, keep $C_1$ & Apply $y'(x_0)$ before continuing \\
A potential you are unsure of & Differentiate it: check $F_x=M$ and $F_y=N$ & Numerical check along one trajectory \\
$M_y=N_x$ but no global potential & Check for a hole: is the origin excluded? & Work on a simply connected subregion \\
$f(x)\dd x=g(y)\dd y$ & Exact for free: $M_y=N_x=0$ & This is separation of variables \\
$\cos(xy)$ or $\eu^{xy}$ throughout & Group $y\dd x+x\dd y$ and set $u=xy$ & Chain rule on top of a known group \\
An implicit answer, $y$ requested & Solve, then pick the branch with the data & Validity ends where $N=0$ \\
\end{recog}
